% ===================================================================
%  TAFEL Nr. 1 — Die Schule. --- Das Gymnasium. --- Das Klassenzimmer.
%
%  Ceci n'est PAS une transcription : c'est une traduction, faite en
%  2026, en allemand d'aujourd'hui. D'ou trois differences avec
%  texto/io et texto/fr, et trois seulement :
%
%   * pas de \nl ni de \cc — la traduction ne suit aucune ligne
%     imprimee, et les lignes de ce fichier ne sont que du confort ;
%   * pas de \begin{VUpage} — elle n'a ni feuillet ni folio, et la
%     page de lecture ne lui en affiche pas ;
%   * le gras se pose sur le TERME seul, ponctuation dehors.
%
%  Tout le reste est identique, et doit l'etre : les cles « %%K » sont
%  celles de texto/io, macro pour macro pour l'apparat, et les renvois
%  (n) visent les MEMES objets de la planche.
%
%  LA SOURCE EST DE 1926, LA COLONNE EST DE 2026. Voir LISEZ-MOI § 5 :
%  le monde decrit reste celui du livret — le Gymnasium, le pion, la
%  diligence — mais la langue est celle d'aujourd'hui. Pas de genitif
%  de chancellerie, pas de « derselbe », pas de « nebst ». L'allemand
%  a ceci de particulier que son « man » impersonnel, lui, n'a pas
%  vieilli d'un jour, la ou l'anglais a du renoncer a son « one » :
%  « dort lernt man Noten lesen » est de 2026 comme de 1926.
%
%  UNE DECISION DE COMPOSITION, ECRITE ICI POUR QU'ON NE LA CHERCHE
%  PAS AILLEURS. L'allemand ecrit « TAFEL Nr. 1 ». « Tabelle » existe,
%  mais c'est le tableau de CHIFFRES ; ce que Delmas vend est une
%  image murale, et c'est « Tafel ». Le mot ne se confond avec aucun
%  autre membre de NUMERO_TAB — aucun ne commence par TAF — et le
%  titre du volume, « Delmas-Hilfstafeln », est en bas de casse. Le
%  cablage a recu « TAFEL », les quatre ordinaux, « Reihe » et
%  « Szene ».
%
%  ET LE MOT DE LA COLONNE EST CELUI DU TABLEAU NOIR, comme en
%  vietnamien : « die Tafel » de l'alinea 2 est l'ardoise murale de la
%  classe, « die Schiefertafel » de l'alinea 12 celle du petit Emil,
%  et « die Lehrtafeln » de l'alinea 2 les planches suspendues. Quatre
%  emplois du meme mot dans la meme page, et le titre est le
%  cinquieme.
%
%  CE QUE LA PAGE MONTRE, ET C'EST LA REPONSE DE TOUTE LA COLONNE :
%  L'IDO COMPOSE SES MOTS COMME L'ALLEMAND.
%
%  Guignon ecrit « desegnochambro » la ou le francais ecrit « classe
%  de dessin » : un mot, deux racines, le determinant devant. C'est la
%  regle allemande, et l'allemand rend le mot sans y toucher —
%  « Zeichensaal ». La page en donne dix autres :
%
%     ido « absorbiva papero »   all. « Löschpapier »   (41)
%     ido « pluvo-tabelo »       all. « Lehrtafel »     (9)
%     ido « skribo-tabuleto »    all. « Schiefertafel » (42)
%     ido « penpozilo »          all. « Federmäppchen » (31)
%     ido « vesteyo-hoki »       all. « Kleiderhaken »  (58)
%     ido « manlaboreyo »        all. « Handarbeitsraum » (69)
%     ido « gimnastik-aparati »  all. « Turngeräte »    (70)
%     ido « skribo-suby »        all. « Schreibunterlage » (53)
%     ido « lern-sako »          all. « Schulranzen »   (57)
%     ido « muziko-pupitro »     all. « Notenständer »  (65)
%
%  Le francais, lui, doit chaque fois trois mots et une preposition :
%  « papier buvard », « tableau mural », « porte-plume », « patères du
%  vestiaire », « salle de travaux manuels ». LE LIVRET EST ECRIT EN
%  IDO, MAIS IL EST PENSE EN ALLEMAND, et cette colonne est la seule
%  des trente ou la traduction soit plus COURTE que l'original.
%
%  ET LA NOTE (*) LE DIT ELLE-MEME, SANS S'EN APERCEVOIR. Guignon
%  enumere sept formes du nom de Jean pour montrer le desordre des
%  langues — « Giovanni, Jean, Johann, Jan, Hans, John, Ivan » — et
%  DEUX DES SEPT SONT ALLEMANDES : Johann et Hans. Aucune autre langue
%  n'en fournit plus d'une. Celle qu'il accuse le plus de multiplier
%  les formes est celle dont il a pris la maniere de faire les mots.
%
%  ENFIN, DEUX MOTS LATINS LA OU L'IDO EN FAIT UN COMPOSE : l'eleve
%  externe et l'eleve interne, « exter-lernanto » et « intern-
%  lernanto » chez Guignon, sont en allemand « der Externe » et « der
%  Interne » — deux adjectifs latins substantives, tels quels. La
%  regle du compose s'arrete a la porte de l'ecole, parce que c'est
%  l'ecole qui a apporte le latin.
%
%  LES NOMS DE PERSONNE SONT ALLEMANDS : Ioannes est Johann, Petrus
%  est Peter, Henrikus est Heinrich, Ludovikus est Ludwig, Georgius
%  est Georg, Iakobus est Jakob. Les noms latins de la note (*) ne
%  bougent pas : c'est d'eux qu'elle parle.
% ===================================================================

%%K t01-apar-0 apar
\VUsaut{2.0mm}
\VUcentre{12.6pt}{120}{ERKLÄRUNGSHEFT}
\VUsaut{3.2mm}
\VUcentre{8.4pt}{160}{ZU DEN}
\VUsaut{2.6mm}
\VUtitre{15.6pt}{0}{Delmas-Hilfstafeln}
\VUsaut{3.4mm}
\VUornamento{9pt}
\VUsaut{5.0mm}
\VUcentre{13.2pt}{90}{ERSTE REIHE}
\VUsaut{3.0mm}
\VUfilet{16mm}
\VUsaut{5.6mm}

%%K t01-apar-1 apar
\VUtitre{15.0pt}{60}{TAFEL Nr. 1}
\VUsaut{1.4mm}
\VUtitre{11.4pt}{0}{Die Schule. --- Das Gymnasium. --- Das Klassenzimmer.}
\VUsaut{2.2mm}
\VUfilet{20mm}
\VUsaut{6.4mm}

%%K t01-tit-1 sub
\VUpk{10.2pt}{40}{Allgemeine Beschreibung.}
\VUsaut{3.0mm}

%%K t01-01-1 p
1. --- Diese Tafel zeigt ein \VUgras{Klassenzimmer} in einem \VUgras{Gymnasium}. In
diesem Klassenzimmer stehen acht \VUgras{Tische} \textsuperscript{(18)} und acht
\VUgras{Bänke} \textsuperscript{(32)}. Auf jeder Bank sitzen drei Schüler. Der
\VUgras{Lehrer} \textsuperscript{(7)} sitzt auf einem \VUgras{Stuhl} \textsuperscript{(8)}
hinter seinem \VUgras{Pult} \textsuperscript{(6)}.

%%K t01-02-1 p
2. --- Links vom Pult steht ein \VUgras{Schrank} mit Glastüren
\textsuperscript{(15)}. Rechts hängt die schwarze \VUgras{Tafel} \textsuperscript{(1)} mit
dem \VUgras{Tafellappen} \textsuperscript{(2)} und der \VUgras{Kreide}
\textsuperscript{(3)}.

%%K t01-02-2 p
Über dem Pult hängen die \VUgras{Lehrtafeln} \textsuperscript{(9, 11, 12)} und eine
\VUgras{Landkarte} \textsuperscript{(10)}.

%%K t01-03-1 p
3. --- Nicht alle Schüler sitzen auf ihrem \VUgras{Platz} \textsuperscript{(17)}.
Johann \textsuperscript{(4)} steht an der Tafel; er hält den Lappen und wischt die
\VUgras{geometrischen Figuren} weg.

%%K t01-03-2 p
Peter steht neben dem Pult; er sammelt die \VUgras{schriftlichen Arbeiten}
\textsuperscript{(5)} ein und bringt sie dem Lehrer.

%%K t01-04-1 p
4. --- Dieser Lehrer unterrichtet \VUgras{moderne Fremdsprachen}. Er bringt den
Schülern bei, fremde Sprachen zu sprechen, zu lesen und zu schreiben
(\VUgras{Deutsch}, \VUgras{Englisch}, \VUgras{Spanisch}, \VUgras{Italienisch},
\VUgras{Russisch} oder \VUgras{Französisch}).

%%K t01-05-1 p
5. --- Das Klassenzimmer ist sehr geräumig und sehr hell. In der Rückwand gibt es
ein breites \VUgras{Fenster}, zwei \VUgras{Oberlichter} \textsuperscript{(78)} und
eine \VUgras{Glastür} zum Lüften und für das Licht.

%%K t01-05-2 p
Kalt ist es hier nicht. In der Mitte steht zum Heizen ein sehr guter
\VUgras{Ofen} \textsuperscript{(46)} mit seinem \VUgras{Ofenrohr} \textsuperscript{(45)}.

%%K t01-05-3 p
An der Wand sind \VUgras{Kleiderhaken} \textsuperscript{(58)} befestigt; daran hängen
die Mützen und Mäntel der Schüler.

%%K t01-tit-2 sub
\VUsaut{5.2mm}
\VUpk{10.2pt}{40}{Der Schulhof.}
\VUsaut{3.6mm}

%%K t01-06-1 p
6. --- Die \VUgras{Uhr} \textsuperscript{(63)} zeigt fünf nach neun.

%%K t01-06-2 p
Die \VUgras{Pause} \textsuperscript{(76)} ist gerade zu Ende, und der Unterricht fängt
wieder an. Die Pause findet auf dem \VUgras{Hof} \textsuperscript{(75)} statt. Wir
sehen ein paar Schüler beim Spielen. Ein \VUgras{Aufseher} \textsuperscript{(74)} hat
sie im Blick.

%%K t01-07-1 p
7. --- Auf dem Hof sehen wir noch andere Leute: den \VUgras{Schulrat}
\textsuperscript{(73)}, den \VUgras{Direktor} \textsuperscript{(71)} und den
\VUgras{Konrektor} \textsuperscript{(72)}.

%%K t01-07-2 p
Der Schulrat prüft die Schulen, der Direktor leitet das Gymnasium, der Konrektor
beaufsichtigt den Unterricht.

%%K t01-07-3 p
Die vierte Person ist der \VUgras{Pförtner} \textsuperscript{(77)}. Er trägt unter dem
Arm ein Heft, in das er die Fehlenden einträgt.

%%K t01-08-1 p
8. --- Hinten auf dem Hof stehen die \VUgras{Turngeräte} \textsuperscript{(70)} und der
\VUgras{Handarbeitsraum} \textsuperscript{(69)}.

%%K t01-08-2 p
Rechts auf der Tafel sehen wir schon die Säle für \VUgras{Musik} und
\VUgras{Zeichnen}.

%%K t01-noto-1 noto
\VUnotes{143.2mm}{%
(*) Bei den \textit{Vornamen} empfiehlt es sich, auch die lateinische Form zu
schreiben. Nur so lassen sich unzählige Sonderformen vermeiden. Wie sollte man
sonst zwischen \textit{Giovanni, Jean,} \textit{Johann, Jan, Hans, John, Ivan}
und so weiter wählen? Sollen wir eigens ein Wörterbuch der Vornamen anlegen?
Nehmen wir aus dem Lateinischen den Nominativ und sagen wir: \VUgras{Ioannes,
Ioanna; Iulius, Iulia; Iakobus; Andreas; Lukas.} (Vollständige Grammatik).%
}

%%K t01-tit-3 sub
\VUpk{10.2pt}{40}{Ausführliche Beschreibung.}
\VUsaut{2.4mm}

%%K t01-tit-4 sub
\VUpk{10.2pt}{40}{Die guten Schüler.}
\VUsaut{3.0mm}

%%K t01-09-1 p
9. --- In der ersten Bank rechts vom Pult sitzen \VUgras{Heinrich} (*),
\VUgras{Stefan} und \VUgras{Karl}.

%%K t01-09-2 p
Heinrich ist sehr aufmerksam und fleißig; er hört dem Lehrer zu. Auf seinem
Tisch liegen ein \VUgras{Heft} \textsuperscript{(28)}, ein \VUgras{Buch}
\textsuperscript{(29)}, ein \VUgras{Wörterbuch} \textsuperscript{(30)} und ein
\VUgras{Federmäppchen} \textsuperscript{(31)}. Sein Buch hat viele \VUgras{Seiten};
jedes Kapitel enthält mehrere \VUgras{Absätze}, und in jeder \VUgras{Zeile} stehen
viele \VUgras{Wörter}.

%%K t01-09-4 p
Sein Nachbar Stefan \textsuperscript{(24)} ist eifrig. Er schreibt sich die Aufgabe
für das nächste Mal ins Heft. Er hat eine schöne \VUgras{Handschrift}. Sein Buch
ist zu. Wir sehen nur den \VUgras{Einband} \textsuperscript{(27)}. Neben ihm liegt
sein \VUgras{Zirkel} \textsuperscript{(26)}.

%%K t01-09-5 p
Karl \textsuperscript{(23)} hält sein Buch in der Hand; er schlägt es auf, um die
Lektion zu wiederholen. Wie jeder Schüler hat er ein \VUgras{Tintenfass}
\textsuperscript{(25)} vor sich. Er ist folgsam.

%%K t01-10-1 p
10. --- Am Tisch daneben sitzen \VUgras{Ludwig}, \VUgras{Anton} und
\VUgras{Paul}. Ludwig sagt seine \VUgras{Lektion} \textsuperscript{(22)} auf und
beantwortet die \VUgras{Fragen} des Lehrers. Anton \textsuperscript{(21)} ist sehr
unaufmerksam; manchmal bestraft ihn der Lehrer sogar. Paul \textsuperscript{(20)} ist
ein guter \VUgras{Kamerad}, freundlich und hilfsbereit. Er ist sehr aufmerksam.

%%K t01-tit-5 sub
\VUsaut{4.2mm}
\VUpk{10.2pt}{40}{Die schlechten Schüler.}
\VUsaut{3.2mm}

%%K t01-11-1 p
11. --- Hinter Heinrich sitzt \VUgras{Georg} \textsuperscript{(38)}. Er ist kein böser
Junge, aber er ist \VUgras{geschwätzig}, unaufmerksam und zappelig. Sein
\VUgras{Leichtsinn} und sein \VUgras{Ungehorsam} bringen \VUgras{Unruhe} in die
Stunde, und oft handelt er sich einen strengen \VUgras{Verweis} ein. Sein
\VUgras{Schmierheft} \textsuperscript{(35)} ist schmutzig, er macht mit seiner
\VUgras{Feder} \textsuperscript{(34)} \VUgras{Tintenkleckse} hinein und braucht
deshalb ständig seinen \VUgras{Radiergummi} \textsuperscript{(36)}; seine
\VUgras{Schulmappe} \textsuperscript{(33)} ist zerrissen, weil er sehr nachlässig ist.
Warum bringt er seinen \VUgras{Federhalter} \textsuperscript{(37)} in die
Fremdsprachenstunde mit?

%%K t01-11-2 p
Sein Nachbar Edmund \textsuperscript{(39)} mag ihn nicht. Er ist zwar ein bisschen
faul, aber er ist still und sitzt ruhig da. Er schwatzt nicht; nur liest er,
statt zuzuhören, lieber die \VUgras{Geschichten} in seinem \VUgras{Lesebuch}.

%%K t01-11-3 p
Der dritte Schüler dieser Bank ist \VUgras{Ludwig} \textsuperscript{(40)}. Er ist
fleißig. Er schreibt auf ein weißes \VUgras{Blatt Papier}. Vor sich hat er sein
\VUgras{Löschpapier} \textsuperscript{(41)} gelegt.

%%K t01-12-1 p
12. --- Die Schüler der zweiten Bank, links vom Lehrer, sind noch sehr jung. Der
erste, der kleine \VUgras{Emil}, kann noch nicht gut schreiben; er bringt seine
\VUgras{Schiefertafel} \textsuperscript{(42)}, seinen \VUgras{Griffel} und seinen
\VUgras{Schwamm} \textsuperscript{(43)} mit.

%%K t01-12-2 p
Neben ihm sitzen \VUgras{August} und \VUgras{Bernhard} \textsuperscript{(44)}. Der
Letzte arbeitet gut, aber seine \VUgras{Kleidung} ist sehr nachlässig.

%%K t01-13-1 p
13. --- Hinter ihnen sitzen drei \VUgras{Faulpelze}. \VUgras{Albert} lernt seine
Lektionen nicht. \VUgras{Jakob} \textsuperscript{(47)} schläft, statt zuzuhören. Seine
\VUgras{Faulheit} bringt ihm \VUgras{Vorwürfe} und sogar \VUgras{Strafen} ein. Und
\VUgras{Christian} \textsuperscript{(48)} schließlich ist viel zu leichtsinnig. Er
\VUgras{benimmt} sich schlecht und gibt sich keine Mühe, sich zu bessern.

%%K t01-14-1 p
14. --- Seine Nachbarn dagegen benehmen sich gut. \VUgras{Ernst}
\textsuperscript{(51)} liebt \VUgras{Ordnung} und \VUgras{Sauberkeit}; er benutzt
immer sein \VUgras{Lineal} \textsuperscript{(50)} und seinen \VUgras{Bleistift}
\textsuperscript{(49)}. Er ist \VUgras{Externer}.

%%K t01-14-2 p
\VUgras{Franz} \textsuperscript{(52)} ist \VUgras{Internatsschüler}; er trägt die
\VUgras{Uniform} des Gymnasiums; er isst und schläft dort. Er ist ein guter
\VUgras{Kamerad}.

%%K t01-14-3 p
Moritz spitzt seinen \VUgras{Bleistift} mit seinem \VUgras{Taschenmesser}
\textsuperscript{(56)}; er ist sehr sorgfältig.

%%K t01-14-4 p
Er bringt alle seine Bücher mit, seine \VUgras{Grammatik} \textsuperscript{(55)}, sein
\VUgras{Notizbuch} \textsuperscript{(54)} und sogar eine \VUgras{Schreibunterlage}
\textsuperscript{(53)}. Sein \VUgras{Schulranzen} \textsuperscript{(57)} steht hinter ihm
auf dem Boden.

%%K t01-14-5 p
Die beiden hinteren Tische im Klassenzimmer haben die \VUgras{guten Schüler}
\textsuperscript{(19)} besetzt.

%%K t01-tit-6 sub
\VUsaut{4.6mm}
\VUpk{10.2pt}{40}{Zeichnen und Musik.}
\VUsaut{3.4mm}

%%K t01-15-1 p
15. --- Im \VUgras{Zeichensaal} hängen schöne \VUgras{Modelle} an der Wand, und von
der Decke hängt eine \VUgras{Lampe} \textsuperscript{(62)} mit einem
\VUgras{Lampenschirm}. Der \VUgras{Lehrer} \textsuperscript{(60)} ist sehr geduldig; er
verbessert die \VUgras{Zeichnungen} \textsuperscript{(61)} der Schüler und bringt
ihnen bei, eine \VUgras{Büste} \textsuperscript{(59)} nach der Natur zu zeichnen.

%%K t01-16-1 p
16. --- Der \VUgras{Musiksaal} sieht sehr interessant aus. Dort lernt man
\VUgras{Noten} lesen, die \VUgras{Töne der Tonleiter} \textsuperscript{(67)} singen,
die auf den \VUgras{Notenlinien} stehen (do, re, mi, fa, sol, la, si\,=\,C. D. E.
F. G. A. B.), die \VUgras{Schlüssel} erkennen und schöne \VUgras{Melodien} singen.
Die Noten stellt man auf den \VUgras{Notenständer} \textsuperscript{(65)}. Einer der
Schüler \VUgras{spielt Klavier} \textsuperscript{(64)}, und der \VUgras{Musiklehrer}
\textsuperscript{(66)} \VUgras{schlägt den Takt} \textsuperscript{(68)}.

%%K t01-17-1 p
17. --- Wir sehen schon eine Ecke des \VUgras{Handarbeitsraums}
\textsuperscript{(69)}, wo die Jungen Holz hobeln und Eisen feilen lernen können.
Nicht jedes Gymnasium hat einen.

%%K t01-tit-7 sub
\VUsaut{5.0mm}
\VUpk{10.2pt}{40}{Der naturwissenschaftliche Saal.}
\VUsaut{3.6mm}

%%K t01-18-1 p
18. --- Der Mathematiklehrer unterrichtet die Schüler in \VUgras{Geometrie},
\VUgras{Arithmetik}, \VUgras{Rechnen} und im \VUgras{metrischen System}. Auf der
Tafel sehen wir die wichtigsten geometrischen Figuren: den \VUgras{Kreis}
\textsuperscript{(a)}, das \VUgras{Quadrat} \textsuperscript{(b)}, das
\VUgras{Dreieck} \textsuperscript{(c)}, die \VUgras{Linie} \textsuperscript{(d)}.

%%K t01-18-2 p
Wir sehen auch eine \VUgras{Rechnung}; es ist keine \VUgras{Addition}, keine
\VUgras{Subtraktion} und keine \VUgras{Multiplikation}: es ist eine
\VUgras{Division} \textsuperscript{(e)}.

%%K t01-19-1 p
19. --- Auf der Tafel des metrischen Systems unterscheiden wir: den
\VUgras{Meter} \textsuperscript{(a)}, die \VUgras{Waage} \textsuperscript{(b)} und ihre
\VUgras{Gewichte}, den \VUgras{Liter} \textsuperscript{(e)}, den \VUgras{Dekaliter}
\textsuperscript{(f)}, den \VUgras{Deziliter} und den \VUgras{Kubikmeter}
\textsuperscript{(d)}.

%%K t01-19-2 p
Die Schüler lernen auch \VUgras{Naturkunde} \textsuperscript{(11)} --- Zoologie
\textsuperscript{(a)}, Botanik \textsuperscript{(b)}, Geologie \textsuperscript{(c)}
--- und \VUgras{Geschichte} \textsuperscript{(12)}.

%%K t01-19-3 p
Im Schrank stehen die Geräte für \VUgras{Physik} \textsuperscript{(15)} und für
\VUgras{Chemie} \textsuperscript{(16)}.

%%K t01-19-4 p
Oben auf dem Schrank stehen: eine \VUgras{Kugel} \textsuperscript{(14)}, ein
\VUgras{Kegel} und ein \VUgras{Globus} \textsuperscript{(13)}.

%%K t01-19-5 p
Über den Tafeln hängt eine \VUgras{Landkarte} mit den fünf Erdteilen:
\VUgras{Europa} \textsuperscript{(g)}, \VUgras{Asien} \textsuperscript{(e)},
\VUgras{Afrika} \textsuperscript{(f)}, \VUgras{Amerika} \textsuperscript{(ab)} und
\VUgras{Ozeanien} \textsuperscript{(h)}, und mit den beiden großen Ozeanen, dem
\VUgras{Pazifik} \textsuperscript{(c)} und dem \VUgras{Atlantik} \textsuperscript{(d)}.
\VUsaut{6.0mm}
\VUfilet{22mm}
